\documentclass[11pt]{article}
\usepackage[margin=1in]{geometry}
\usepackage{amsmath,amssymb,amsthm}
\usepackage[T1]{fontenc}
\title{A countability obstruction to a full interval of graph-invariant gaps\\
  A disproof of one assertion of conjecture 00000003480}
\author{earthking11}
\date{}

\begin{document}
\maketitle

Conjecture 00000003480 states that the difference between the fractional
chromatic number and a fractional Hall bound takes every real value from
$0$ to $1$, with a Mycielski-type family realizing the values. For finite
unweighted simple graphs, such a full interval is impossible for any pair of
real-valued graph invariants.

For each $n\geq 0$ there are at most $2^{\binom n2}$ simple graphs on a
labeled $n$-element vertex set. Hence the disjoint union of the graph sets
over all $n$ is countable. Every finite simple graph is isomorphic to one
in this union. If $A$ and $B$ are any real-valued graph invariants, their
difference has at most countably many values:
\[
  \bigl\{A(G)-B(G):G\text{ a finite simple graph}\bigr\}
  \quad\text{is countable}.
\]
In contrast, $[0,1]$ is uncountable. For example, the injective map
\[
  x\longmapsto\frac{\arctan(x)+\pi/2}{\pi}
\]
sends $\mathbb R$ into $(0,1)$. Therefore
\[
  [0,1]\not\subseteq
  \bigl\{\chi_f(G)-H_f(G):G\text{ a finite simple graph}\bigr\},
\]
where $\chi_f$ and $H_f$ denote the two invariants in the conjecture.
The asserted continuous spectrum fails, regardless of the definitions of
these invariants or the behavior of the endpoint $1$.

\paragraph{Formalization.}
The Lean file encodes the class of all finite labeled simple graphs as
\texttt{$\Sigma$ n : $\mathbb N$, SimpleGraph (Fin n)}. It proves this class
countable and the real interval $[0,1]$ uncountable. The theorem
\texttt{TLMC3480.no\_full\_interval\_of\_difference} states that the
difference of any two real-valued functions on this graph class cannot take
every value in $[0,1]$. The arbitrary-function statement applies to the
named fractional graph invariants without using their LP formulas. Its
axiom audit contains only \texttt{propext}, \texttt{Classical.choice}, and
\texttt{Quot.sound}.

\paragraph{Scope.}
The argument addresses finite unweighted simple graphs, the standard domain
of the described Mycielski construction. A differently stated conjecture
allowing continuously weighted or otherwise uncountably parametrized graph
objects is outside its scope.

\end{document}
